\documentclass[11pt,letterpaper]{article}

\usepackage[margin=0.82in,headheight=15pt,footskip=28pt]{geometry}
\usepackage{amsthm}
\usepackage{newtxtext,newtxmath}
\usepackage[T1]{fontenc}
\usepackage{microtype}
\usepackage{mathtools,bm}
\usepackage{booktabs,longtable,tabularx,array}
\usepackage{enumitem}
\usepackage{xcolor}
\usepackage[most]{tcolorbox}
\usepackage{titlesec}
\usepackage{fancyhdr}
\usepackage{hyperref}
\usepackage{lastpage}

\definecolor{onyx}{HTML}{17191C}
\definecolor{graphite}{HTML}{4B4F55}
\definecolor{rulegray}{HTML}{B8BBC0}
\definecolor{papergray}{HTML}{F3F3F1}
\definecolor{gold}{HTML}{8B6B2D}
\definecolor{certgreen}{HTML}{315B45}
\definecolor{declareblue}{HTML}{31556F}

\hypersetup{
  colorlinks=true,
  linkcolor=onyx,
  urlcolor=declareblue,
  pdftitle={Clear_Onyx: Computational Kernel Briefing},
  pdfauthor={AVC},
  pdfsubject={Exact R2 operator algebra and computational architecture},
  pdfkeywords={Clear_Onyx, matrix algebra, Clifford algebra, Fibonacci operator, normal form, computational kernel}
}

\pagestyle{fancy}
\fancyhf{}
\fancyhead[L]{\small\textsc{Clear\_Onyx}}
\fancyhead[R]{\small Computational Kernel Briefing}
\fancyfoot[L]{\small AVC -- 5 August 2026}
\fancyfoot[R]{\small \thepage\ of \pageref{LastPage}}
\renewcommand{\headrulewidth}{0.35pt}
\renewcommand{\footrulewidth}{0pt}

\titleformat{\section}
  {\Large\bfseries\color{onyx}}
  {\thesection}{0.7em}{}
  [\vspace{0.2em}\color{rulegray}\titlerule]
\titleformat{\subsection}
  {\large\bfseries\color{onyx}}
  {\thesubsection}{0.65em}{}
\titleformat{\subsubsection}
  {\normalsize\bfseries\color{graphite}}
  {\thesubsubsection}{0.6em}{}
\titlespacing*{\section}{0pt}{2.2ex plus .4ex minus .2ex}{1.2ex}
\titlespacing*{\subsection}{0pt}{1.8ex plus .3ex minus .2ex}{0.8ex}

\setlist[itemize]{leftmargin=1.4em,itemsep=0.25em,topsep=0.35em}
\setlist[enumerate]{leftmargin=1.55em,itemsep=0.25em,topsep=0.35em}
\setlength{\parindent}{0pt}
\setlength{\parskip}{0.55em}
\renewcommand{\arraystretch}{1.18}

\newtheoremstyle{clearonyx}
  {0.7em}{0.7em}{\itshape}{}% space above/below, body font, indent
  {\bfseries\color{onyx}}{.}{0.5em}{}
\theoremstyle{clearonyx}
\newtheorem{proposition}{Proposition}[section]
\newtheorem{corollary}[proposition]{Corollary}
\newtheorem{definition}[proposition]{Definition}
\newtheorem{remark}[proposition]{Remark}

\newcommand{\R}{\mathbb{R}}
\newcommand{\Qfield}{\mathbb{Q}}
\newcommand{\Z}{\mathbb{Z}}
\newcommand{\N}{\mathbb{N}}
\newcommand{\NF}{\operatorname{NF}}
\newcommand{\Eval}{\operatorname{Eval}}
\newcommand{\Fix}{\operatorname{Fix}}
\newcommand{\tr}{\operatorname{tr}}
\newcommand{\diag}{\mathrm{diag}}
\newcommand{\Span}{\operatorname{span}}
\newcommand{\Cl}{\mathrm{Cl}}
\newcommand{\Append}{\mathbin{\frown}}
\newcommand{\statuscert}{\textcolor{certgreen}{\textbf{CERT}}}
\newcommand{\statusderived}{\textcolor{gold}{\textbf{DERIVED}}}
\newcommand{\statusdeclared}{\textcolor{declareblue}{\textbf{DECLARED}}}

\newtcolorbox{executivebox}{
  colback=papergray,
  colframe=rulegray,
  boxrule=0.5pt,
  arc=1.2pt,
  left=8pt,right=8pt,top=7pt,bottom=7pt
}
\newtcolorbox{boundarybox}{
  colback=white,
  colframe=gold,
  boxrule=0.65pt,
  arc=1.2pt,
  left=8pt,right=8pt,top=7pt,bottom=7pt,
  title={\bfseries Evidence boundary},
  coltitle=onyx,
  fonttitle=\normalsize
}
\newtcolorbox{resultbox}{
  colback=papergray,
  colframe=onyx,
  boxrule=0.7pt,
  arc=1.2pt,
  left=9pt,right=9pt,top=8pt,bottom=8pt
}

\begin{document}

\begin{titlepage}
\thispagestyle{empty}
\vspace*{0.12\textheight}

{\fontsize{34}{38}\selectfont\bfseries\color{onyx} CLEAR\_ONYX\par}
\vspace{0.75em}
{\fontsize{22}{27}\selectfont\bfseries Computational Kernel Briefing\par}
\vspace{0.75em}
{\large Exact $\R^2$ Operator Algebra, Normal Forms, Phase/Lift Addressing, and Typed Composition\par}

\vspace{2.4em}
\color{gold}\rule{0.56\textwidth}{1.1pt}\color{black}
\vspace{2.2em}

\begin{executivebox}
\textbf{Central result.}
The kernel generated by the Fibonacci operator $Q$ and the quarter-turn operator $C_1$ is not merely compact.  It spans the full real matrix algebra $M_2(\R)$, gives every $2\times2$ real operator a unique four-coordinate normal form, reduces arbitrary finite operator words exactly, and separates cyclic visible phase from unbounded retained occurrence.  The result is a small exact engine with infinite compositional syntax and finite canonical semantics.
\end{executivebox}

\vfill
{\large AVC\par}
{\large 5 August 2026\par}
\vspace{0.6em}
{\small Technical briefing -- print edition\par}
\end{titlepage}

\tableofcontents
\clearpage

\section{Purpose and scope}

Once the following algebra is organized as a computational engine, it provides substantial freedom to compose while keeping every composition inside a small exact constraint system.  Arbitrarily long words in the generating operators can be preserved as provenance, evaluated exactly, and reduced to a unique canonical state.  The same kernel can therefore support many lawful rendering, semantic, or physical adapters without allowing those adapters to rewrite the certified mathematics.

\begin{boundarybox}
This note separates three evidence levels.

\begin{itemize}
  \item \statuscert\ denotes identities proved directly from the displayed matrices.
  \item \statusderived\ denotes consequences proved from certified identities together with an explicitly named design choice.
  \item \statusdeclared\ denotes architecture selected for the platform, such as the $137$-division angular register and the eight-channel carrier.  A declared register may be useful and exact as a data type without being forced by the $2\times2$ kernel itself.
\end{itemize}

The note certifies an operator algebra and a computational organization.  It does not, by itself, certify any particular physical interpretation attached through an adapter.
\end{boundarybox}

\section{Exact mathematical kernel}

\subsection{Base matrices}

The kernel begins with four explicitly represented matrices acting on $\R^2$:

\begin{equation}
Q=
\begin{pmatrix}
1&1\\
1&0
\end{pmatrix},
\qquad
\Omega=
\begin{pmatrix}
0&1\\
-1&0
\end{pmatrix},
\qquad
I=
\begin{pmatrix}
1&0\\
0&1
\end{pmatrix},
\qquad
C_1=
\begin{pmatrix}
0&-1\\
1&0
\end{pmatrix}=-\Omega.
\end{equation}

The matrix $Q$ is symmetric and has characteristic polynomial

\begin{equation}
\chi_Q(\lambda)=\lambda^2-\lambda-1.
\end{equation}

Its eigenvalues are the golden conjugates

\begin{equation}
\Phi=\frac{1+\sqrt5}{2},
\qquad
\Psi=\frac{1-\sqrt5}{2}=-\Phi^{-1},
\end{equation}

with

\begin{equation}
\Phi+\Psi=1,
\qquad
\Phi\Psi=-1,
\qquad
\Phi-\Psi=\sqrt5.
\end{equation}

\subsection{Generated quadratic operator}

Define the order-sensitive operator

\begin{equation}
G=Q^{T}\Omega-\Omega Q.
\end{equation}

Since $Q^T=Q$, direct multiplication gives

\begin{equation}
G=
\begin{pmatrix}
-2&1\\
1&2
\end{pmatrix}.
\end{equation}

The normalized operator is

\begin{equation}
H_{\sqrt5}=\frac{G}{\sqrt5}.
\end{equation}

\subsection{Scalar and register data}

The exact scalar layer uses

\begin{equation}
\sqrt5,
\qquad
\Phi=\frac{1+\sqrt5}{2},
\qquad
\tau=2\pi.
\end{equation}

The platform architecture additionally uses the following typed registers:

\begin{center}
\begin{tabularx}{0.92\textwidth}{@{}>{\raggedright\arraybackslash}p{0.28\textwidth}X>{\raggedright\arraybackslash}p{0.18\textwidth}@{}}
\toprule
\textbf{Register} & \textbf{Role} & \textbf{Status}\\
\midrule
$137$ angular divisions & Exact cyclic angular address $a\in\Z/137\Z$ & \statusdeclared\\
$8$ addressed states & Channel carrier for relational payloads & \statusdeclared\\
$28$ pair relations & Complete unordered pair set among eight channels & \statusderived\\
$4$ visible phase positions & Quarter-turn phase class $p\in\Z/4\Z$ & \statuscert\\
$c\in\N$ & Retained completed-lift count & \statusdeclared\\
$p\in\{0,1,2,3\}$ & Visible phase representative & \statuscert\\
$n=4c+p$ & Total occurrence address & \statusderived\\
\bottomrule
\end{tabularx}
\end{center}

\section{Certificate ledger}

All identities in this section follow by exact matrix multiplication.

\begin{longtable}{@{}p{0.34\textwidth}p{0.55\textwidth}p{0.09\textwidth}@{}}
\toprule
\textbf{Identity} & \textbf{Operational consequence} & \textbf{Status}\\
\midrule
\endfirsthead
\toprule
\textbf{Identity} & \textbf{Operational consequence} & \textbf{Status}\\
\midrule
\endhead
$Q^2=Q+I$ & Fibonacci lift closes quadratically; every power of $Q$ reduces to the span of $I$ and $Q$. & \statuscert\\
$Q^T\Omega Q=-\Omega$ & One $Q$-lift reverses the oriented area form; two lifts preserve it. & \statuscert\\
$Q^TGQ=-G$ & One $Q$-lift reverses the sign of the $G$-audit form; two lifts preserve it. & \statuscert\\
$G^2=5I$ & The generated audit has exact scale $\sqrt5$. & \statuscert\\
$H_{\sqrt5}^2=I$ & The normalized audit is an involution. & \statuscert\\
$H_{\sqrt5}^{T}H_{\sqrt5}=I$ & The normalized involution is Euclidean orthogonal. & \statuscert\\
$C_1^2=-I$, $C_1^4=I$ & $C_1$ is an exact quarter-turn with four visible phase positions. & \statuscert\\
$\binom82=28$ & Eight declared channels have exactly twenty-eight unordered pair relations. & \statusderived\\
$\Delta_G=0$ on $y/x=1/\Phi$ and $y/x=-\Phi$ & The quadratic audit vanishes exactly on the two golden eigendirections of $Q$. & \statuscert\\
$\Phi^{-c}Q^c$ remains bounded & Normalized lift can continue for unbounded $c$ without exponential state magnitude. & \statuscert\\
$n=4c+p$ & Cyclic phase and retained occurrence have a unique quotient-remainder decomposition. & \statusderived\\
\bottomrule
\end{longtable}

Two further invariants are useful:

\begin{equation}
\tr(G)=0,
\qquad
\det(G)=-5,
\qquad
\det(H_{\sqrt5})=-1.
\end{equation}

Thus $H_{\sqrt5}$ is a symmetric orthogonal reflection with eigenvalues $+1$ and $-1$.

\section{Algebraic completeness and unique normal form}

The reason this kernel is useful is stronger than its small size: it supplies the full operator algebra on a two-dimensional real state space.

\begin{proposition}[Basis theorem]
The ordered set
\[
\mathcal B=(I,Q,\Omega,G)
\]
is a basis of $M_2(\R)$.
\end{proposition}

\begin{proof}
Vectorize each matrix in row-major order and place the vectors as columns:

\begin{equation}
B=
\begin{pmatrix}
1&1&0&-2\\
0&1&1&1\\
0&1&-1&1\\
1&0&0&2
\end{pmatrix}.
\end{equation}

Direct evaluation gives

\begin{equation}
\det(B)=-10\neq0.
\end{equation}

Therefore the four matrices are linearly independent.  Since $\dim_{\R}M_2(\R)=4$, they form a basis.
\end{proof}

\begin{corollary}[Canonical coordinates]
Every matrix
\[
A=\begin{pmatrix}a&b\\c&d\end{pmatrix}\in M_2(\R)
\]
has one and only one normal form
\[
A=x_0I+x_1Q+x_2\Omega+x_3G.
\]
Its coordinates are
\begin{align}
 x_0&=\frac{2a-b-c+3d}{5},\\
 x_1&=\frac{a+2b+2c-d}{5},\\
 x_2&=\frac{b-c}{2},\\
 x_3&=\frac{-2a+b+c+2d}{10}.
\end{align}
\end{corollary}

This gives every evaluated construction an exact four-coordinate representation.  A long operator word and a short operator word are semantically identical precisely when their four normal-form coordinates agree.

\begin{resultbox}
\begin{equation*}
\boxed{\Span_{\R}\{I,Q,\Omega,G\}=M_2(\R)}
\end{equation*}

The kernel is a complete four-dimensional real matrix algebra acting on $\R^2$, not a special-purpose animation rule.
\end{resultbox}

\section{Two generators and exact word reduction}

\subsection{Minimal generating alphabet}

Because $C_1=-\Omega$ and $Q^T=Q$,

\begin{equation}
G=Q\Omega-\Omega Q=C_1Q-QC_1=[C_1,Q].
\end{equation}

Thus the full algebra is generated by the two-letter alphabet

\begin{equation}
\boxed{\{Q,C_1\}}.
\end{equation}

The relation

\begin{equation}
Q^2=Q+I
\end{equation}

makes $Q$ the Fibonacci growth or lift operator.  The relations

\begin{equation}
C_1^2=-I,
\qquad
C_1^4=I
\end{equation}

make $C_1$ the quarter-turn phase operator.  Their commutator

\begin{equation}
G=C_1Q-QC_1
\end{equation}

measures the exact difference between applying turn then growth and applying growth then turn.

\begin{resultbox}
\begin{equation*}
\boxed{G=\text{order-sensitive composition audit}.}
\end{equation*}
\end{resultbox}

\subsection{Finite reduction table}

Let

\begin{equation}
w=g_1g_2\cdots g_m,
\qquad
g_j\in\{Q,C_1\}
\end{equation}

be any finite operator word.  Its evaluated matrix reduces uniquely to

\begin{equation}
\NF(w)=x_0I+x_1Q+x_2\Omega+x_3G.
\end{equation}

All products reduce through the finite table

\begin{align}
Q^2&=Q+I,
&\Omega^2&=-I,
&G^2&=5I,\\
Q\Omega&=\frac{\Omega+G}{2},
&\Omega Q&=\frac{\Omega-G}{2},\\
QG&=\frac{G+5\Omega}{2},
&GQ&=\frac{G-5\Omega}{2},\\
\Omega G&=2Q-I,
&G\Omega&=I-2Q.
\end{align}

Therefore the engine has

\begin{equation}
\boxed{\text{infinite compositional syntax}\; +\; \text{finite exact semantics}.}
\end{equation}

The unreduced word can be retained as provenance while the four coefficients provide exact comparison, caching, testing, rendering input, and error localization.

\subsection{Exact coordinate multiplication}

If

\begin{equation}
A=x_0I+x_1Q+x_2\Omega+x_3G,
\qquad
B=y_0I+y_1Q+y_2\Omega+y_3G,
\end{equation}

then $AB=z_0I+z_1Q+z_2\Omega+z_3G$, where

\begin{align}
 z_0={}&x_0y_0+x_1y_1-x_2y_2-x_2y_3+x_3y_2+5x_3y_3,\\
 z_1={}&x_0y_1+x_1y_0+x_1y_1+2x_2y_3-2x_3y_2,\\
 z_2={}&x_0y_2+\frac{x_1y_2}{2}+\frac{5x_1y_3}{2}
       +x_2y_0+\frac{x_2y_1}{2}-\frac{5x_3y_1}{2},\\
 z_3={}&x_0y_3+\frac{x_1y_2}{2}+\frac{x_1y_3}{2}
       -\frac{x_2y_1}{2}+x_3y_0+\frac{x_3y_1}{2}.
\end{align}

This bilinear law is a direct implementation interface: multiplication can occur entirely in normal-form coordinates without reconstructing an intermediate symbolic expression.

\section{Normalized Clifford core}

Set

\begin{equation}
H_{\sqrt5}=\frac{G}{\sqrt5}.
\end{equation}

Then

\begin{equation}
H_{\sqrt5}^2=I,
\qquad
\Omega^2=-I,
\qquad
H_{\sqrt5}\Omega=-\Omega H_{\sqrt5}.
\end{equation}

These are the defining relations of the real Clifford algebra of one positive and one negative generator:

\begin{equation}
\Cl_{1,1}(\R)\cong M_2(\R).
\end{equation}

A normalized operational basis is therefore

\begin{equation}
\{I,H_{\sqrt5},\Omega,\Omega H_{\sqrt5}\}.
\end{equation}

The original lift operator is recoverable from the normalized core.  Since $\Omega G=2Q-I$,

\begin{equation}
Q=\frac{I+\Omega G}{2}
 =\frac{I+\sqrt5\,\Omega H_{\sqrt5}}{2}
 =\frac{I-\sqrt5\,C_1H_{\sqrt5}}{2}.
\end{equation}

Thus growth, quarter-turn phase, reflection, and the mixed direction are not separate modules.  They are coordinated offices in one exact algebra.

\section{The golden seam as audit boundary}

For $v=(x,y)^T$, define the quadratic audit

\begin{equation}
\Delta_G(v)=v^TGv.
\end{equation}

Direct substitution gives

\begin{equation}
\Delta_G(x,y)=-2x^2+2xy+2y^2.
\end{equation}

Using $\Phi-\Phi^{-1}=1$, this factors exactly as

\begin{equation}
\boxed{
\Delta_G(x,y)
=2\left(y-\frac{x}{\Phi}\right)\left(y+\Phi x\right).
}
\end{equation}

Hence the null directions are

\begin{equation}
\frac{y}{x}=\frac1\Phi,
\qquad
\frac{y}{x}=-\Phi.
\end{equation}

These are precisely the eigendirections of $Q$:

\begin{equation}
v_\Phi=(\Phi,1)^T,
\qquad
v_\Psi=(\Psi,1)^T,
\end{equation}

with

\begin{equation}
Qv_\Phi=\Phi v_\Phi,
\qquad
Qv_\Psi=\Psi v_\Psi.
\end{equation}

The normalized involution exchanges the two golden directions projectively:

\begin{equation}
H_{\sqrt5}v_\Phi=\Phi v_\Psi,
\qquad
H_{\sqrt5}v_\Psi=\Phi^{-1}v_\Phi.
\end{equation}

Therefore

\begin{equation}
\boxed{
H_{\sqrt5}=\text{the exact involution exchanging the two golden growth directions}.
}
\end{equation}

The identities $Q^TGQ=-G$ and $(Q^2)^TGQ^2=G$ add a useful parity interpretation: a single lift crosses the sign of the quadratic audit, while a completed pair of lifts preserves it.

\section{Bounded normalized Fibonacci evolution}

The quadratic relation for $Q$ yields the exact Fibonacci power law

\begin{equation}
Q^c=F_cQ+F_{c-1}I,
\qquad c\ge1,
\end{equation}

where $F_0=0$, $F_1=1$, and $F_{c+1}=F_c+F_{c-1}$.

Normalize by the dominant eigenvalue:

\begin{equation}
\widehat Q_c=\Phi^{-c}Q^c.
\end{equation}

The spectral projectors are

\begin{equation}
P_\Phi=\frac{Q-\Psi I}{\sqrt5},
\qquad
P_\Psi=\frac{\Phi I-Q}{\sqrt5}.
\end{equation}

They satisfy

\begin{equation}
P_\Phi+P_\Psi=I,
\qquad
P_\Phi P_\Psi=0,
\qquad
Q=\Phi P_\Phi+\Psi P_\Psi.
\end{equation}

Consequently,

\begin{equation}
\boxed{
\Phi^{-c}Q^c
=P_\Phi+\left(-\Phi^{-2}\right)^cP_\Psi.
}
\end{equation}

Because $|\Phi^{-2}|<1$, the normalized sequence is bounded and converges exponentially to $P_\Phi$.

Computationally, this supplies:

\begin{itemize}
  \item an unbounded lift count $c$;
  \item bounded normalized operator magnitude;
  \item an exact retained coefficient for the conjugate mode;
  \item a stable asymptotic state;
  \item exact discrete kernel states that do not need to be inferred from interpolation between unrelated rendered endpoints.
\end{itemize}

The occurrence count can grow without bound while the normalized state remains controlled.

\section{Operational offices}

The kernel becomes an instrument when each object has a single explicit computational office.

\begin{longtable}{@{}p{0.28\textwidth}p{0.62\textwidth}p{0.08\textwidth}@{}}
\toprule
\textbf{Object} & \textbf{Operational office} & \textbf{Status}\\
\midrule
\endfirsthead
\toprule
\textbf{Object} & \textbf{Operational office} & \textbf{Status}\\
\midrule
\endhead
$I$ & Identity retention: preserve the current state. & \statuscert\\
$C_1=-\Omega$ & Positive quarter-turn phase: open the orthogonal direction. & \statuscert\\
$\Omega$ & Oriented area form and signed rotational carrier. & \statuscert\\
$Q$ & Fibonacci lift: generate the next exact growth state. & \statuscert\\
$G=[C_1,Q]$ & Detect the difference made by operator order. & \statuscert\\
$H_{\sqrt5}=G/\sqrt5$ & Normalize the order difference into an exact orthogonal involution. & \statuscert\\
$\Delta_G(v)$ & Classify a direction relative to the two golden null seams. & \statuscert\\
$\Phi^{-c}$ & Normalize accumulated lift while preserving its exact spectral form. & \statuscert\\
$p\in\Z/4\Z$ & Visible phase position. & \statuscert\\
$c\in\N$ & Retained count of completed four-phase lifts. & \statusdeclared\\
$n=4c+p$ & Unambiguous total occurrence address. & \statusderived\\
$a\in\Z/137\Z$ & Exact angular scope address. & \statusdeclared\\
Eight states & Addressed channel carrier. & \statusdeclared\\
Twenty-eight pairs & Complete relational seam set. & \statusderived\\
Append-only $R_n$ & Preserve occurrence even when visible phase returns. & \statusdeclared\\
\bottomrule
\end{longtable}

The phrase ``every $+1$'' has a literal implementation meaning:

\begin{equation}
\boxed{+1=\text{apply one typed transition and retain its occurrence}.}
\end{equation}

It does not mean merely draw one more visual frame.

\section{Phase return and retained occurrence}

Let $n\in\N$.  Euclidean division by four gives the unique decomposition

\begin{equation}
n=4c+p,
\qquad
0\le p<4,
\end{equation}

with

\begin{equation}
p=n\bmod4,
\qquad
c=\left\lfloor\frac n4\right\rfloor.
\end{equation}

Thus visible phase and retained history are separated:

\begin{equation}
p_{n+4}=p_n,
\qquad
c_{n+4}=c_n+1.
\end{equation}

The exact computational statement is

\begin{equation}
\boxed{\text{phase closes; occurrence continues}.}
\end{equation}

The angular register is separately typed as

\begin{equation}
a\in\Z/137\Z.
\end{equation}

Since $\gcd(4,137)=1$, the Chinese remainder theorem gives

\begin{equation}
\Z/4\Z\times\Z/137\Z\cong\Z/548\Z.
\end{equation}

The finite $548$-address register joins visible phase and angular address exactly.  The unbounded count $c$ remains outside it.  The system may therefore revisit a finite visible address without revisiting the same occurrence.

For exact angular storage, retain the integer address and interpret

\begin{equation}
\theta_a=\frac{a}{137}\tau
\end{equation}

symbolically.  Floating-point sine, cosine, color conversion, and pixels belong at the rendering boundary rather than in the exact address layer.

\section{Eight channels and twenty-eight relational seams}

Let the declared channel carrier be

\begin{equation}
A_8\cong\R^8.
\end{equation}

Its complete pairwise relation space is the second exterior power

\begin{equation}
\Lambda^2(A_8^*),
\end{equation}

with dimension

\begin{equation}
\dim\Lambda^2(A_8^*)=\binom82=28.
\end{equation}

The eight channels are a design choice.  Once they are chosen, the twenty-eight unordered pair relations are forced.

If each channel carries a phase, a common phase is relationally invisible.  The relative-phase carrier is

\begin{equation}
U(1)^8/U(1)_{\diag}\cong U(1)^7.
\end{equation}

The architecture therefore keeps distinct:

\begin{itemize}
  \item eight addressed channel states;
  \item seven independent relative-phase degrees of freedom;
  \item twenty-eight complete pairwise reads;
  \item four visible phase positions;
  \item $137$ angular addresses;
  \item one unbounded retained lift count.
\end{itemize}

These axes are typed separately rather than collapsed into one opaque index.

\section{Exact computational state and transition contract}

A minimal state record may be written

\begin{equation}
X_n=\left(A_n,p_n,c_n,a_n,u_n,R_n\right),
\end{equation}

where

\begin{equation}
A_n=x_{0,n}I+x_{1,n}Q+x_{2,n}\Omega+x_{3,n}G
\end{equation}

is the exact operator normal form,

\begin{equation}
p_n\in\Z/4\Z,
\qquad
c_n\in\N,
\qquad
a_n\in\Z/137\Z,
\qquad
u_n\in K^8,
\end{equation}

and $K$ is the scalar or phase domain declared by the active adapter.  The occurrence ledger is append-only:

\begin{equation}
R_{n+1}=R_n\Append o_{n+1}.
\end{equation}

Every transition should retain at least:

\begin{enumerate}
  \item the generator applied;
  \item the unreduced operator word;
  \item the exact reduced normal form;
  \item the input and output addresses;
  \item the audit result;
  \item any rendered approximation produced at the boundary.
\end{enumerate}

This record allows the platform to identify where a fracture entered rather than merely displaying a wrong final result.

A typed transition may be represented schematically as

\begin{equation}
T_g:\;X_n\longmapsto X_{n+1},
\qquad
g\in\{Q,C_1\},
\end{equation}

subject to the invariants

\begin{equation}
A_{n+1}=\NF(A_ng),
\qquad
p_{n+1}=p_n+1\pmod4,
\qquad
R_{n+1}=R_n\Append o_{n+1},
\end{equation}

with adapter-specific updates to $a_n$ and $u_n$ declared separately.

\section{Where free composition lives}

The certified identities remain fixed.  Compositional freedom lies in choosing:

\begin{itemize}
  \item which legal generator word to execute;
  \item which normal-form coordinate or invariant to observe;
  \item which channel or seam to read;
  \item which angular address to expose;
  \item which semantic or physical adapter interprets the state;
  \item which visual grammar renders it;
  \item which typed transition policy selects the next legal word.
\end{itemize}

A color loom, symmetry viewer, phase clock, grammatical construction, or physical simulation may therefore be implemented as a distinct adapter around the same core.  Each adapter must declare what its coordinates mean and must not promote visual resemblance into an algebraic or physical certificate.

The architecture can be summarized as

\begin{equation}
\boxed{
\begin{aligned}
\text{operator words}
&\xrightarrow{\Eval}
M_2\!\left(\Qfield(\sqrt5)\right)
\xrightarrow{\NF}
(x_0,x_1,x_2,x_3)\\
&\xrightarrow{\text{typed address}}
(p,c,a,u,\text{seam},R)
\xrightarrow{\text{adapter}}
\text{visible or physical realization}.
\end{aligned}}
\end{equation}

For words in the unnormalized generators $Q$ and $C_1$, evaluation already lies in $M_2(\Z)$ and normal-form coordinates lie in $\Qfield$.  The extension $\Qfield(\sqrt5)$ enters when the normalized involution $H_{\sqrt5}$ or golden spectral projectors are used.

\begin{resultbox}
The core provides the five structural properties required of a reusable generative engine:

\begin{equation*}
\boxed{
\text{exact syntax}
+\text{canonical reduction}
+\text{bounded normalized evolution}
+\text{unbounded memory}
+\text{typed free adapters}.
}
\end{equation*}
\end{resultbox}

The kernel is not the finished composition.  It is the small exact machine that permits many compositions without allowing any one of them to quietly rewrite the rules.

\clearpage
\section*{Implementation reference card}
\addcontentsline{toc}{section}{Implementation reference card}

\begin{executivebox}
\textbf{Generators}
\begin{equation*}
Q=\begin{pmatrix}1&1\\1&0\end{pmatrix},
\qquad
C_1=\begin{pmatrix}0&-1\\1&0\end{pmatrix},
\qquad
\Omega=-C_1.
\end{equation*}

\textbf{Generated operators}
\begin{equation*}
G=[C_1,Q]=\begin{pmatrix}-2&1\\1&2\end{pmatrix},
\qquad
H_{\sqrt5}=G/\sqrt5.
\end{equation*}

\textbf{Core relations}
\begin{gather*}
Q^2=Q+I,
\qquad
C_1^2=-I,
\qquad
G^2=5I,
\qquad
H_{\sqrt5}^2=I,\\
Q^T\Omega Q=-\Omega,
\qquad
Q^TGQ=-G,
\qquad
H_{\sqrt5}\Omega=-\Omega H_{\sqrt5}.
\end{gather*}

\textbf{Normal form}
\begin{equation*}
A=x_0I+x_1Q+x_2\Omega+x_3G.
\end{equation*}

\textbf{Golden audit}
\begin{equation*}
\Delta_G(x,y)=2\left(y-\frac{x}{\Phi}\right)(y+\Phi x),
\qquad
\Phi=\frac{1+\sqrt5}{2}.
\end{equation*}

\textbf{Normalized lift}
\begin{equation*}
\Phi^{-c}Q^c=P_\Phi+(-\Phi^{-2})^cP_\Psi.
\end{equation*}

\textbf{Address}
\begin{equation*}
n=4c+p,
\qquad
p=n\bmod4,
\qquad
c=\lfloor n/4\rfloor,
\qquad
a\in\Z/137\Z.
\end{equation*}

\textbf{Relational carrier}
\begin{equation*}
\dim\Lambda^2((\R^8)^*)=\binom82=28,
\qquad
U(1)^8/U(1)_{\diag}\cong U(1)^7.
\end{equation*}

\textbf{State}
\begin{equation*}
X_n=(A_n,p_n,c_n,a_n,u_n,R_n),
\qquad
R_{n+1}=R_n\Append o_{n+1}.
\end{equation*}
\end{executivebox}

\vfill
\begin{center}
\small\color{graphite}
Clear\_Onyx computational kernel briefing -- AVC -- 5 August 2026
\end{center}

\end{document}
